\section{GeoToolEval}
\label{sec:benchmark}
GeoToolEval evaluates tool use and multistep orchestration through records containing a natural-language question, input resources, a reference answer and answer type, and reference tool calls. Single-step tasks assess data location, parameter determination, and basic invocation; multistep tasks examine operation composition, dependencies, and use of intermediate results. Tasks span ten research themes, with theme sizes determined by available resources and screening outcomes.

\subsection{Sources and Construction}
Tasks are either selected from existing question-answer resources or derived from scientific data. Geoscientific big-data analysis includes 40 selected GeoQuestions1089 tasks involving physical geographic objects and spatial operations. Remote sensing includes 49 of Earth-Bench's 248 tasks, retaining source questions, choices, and answers while adapting calls to our interfaces. Critical-mineral and climate-change tasks are selected from \nolinkurl{habitable_earth_dataset}; underlying resources include MRDS, paleoclimate reconstructions, GISTEMP, and SSP warming summaries. Data-derived tasks specify identifiable objects, variables, and computational relationships, with answers determined through processing and calculation. Table~\ref{tab:benchmark} lists the data sources.

Reference calls map source operational requirements to registered tools and record parameters and input/output information. Existing tasks retain source identifiers or resource locators; derived tasks retain data-location fields, derivation records, and computational scripts. Quality control checks traceability, record completeness, consistency of reference invocations, and sample execution. Sample validation is recorded separately from full-benchmark answer evaluation.

% DOCX paragraphs 119--119, table 1. All counts are copied from the source.
\begin{table*}[t]
  \caption{GeoToolEval task sources and reference tool-invocation chain lengths. Counts refer to evaluation records.}
  \label{tab:benchmark}
  \centering
  \begin{tabularx}{\textwidth}{@{}p{0.25\textwidth}Xrrrr@{}}
    \toprule
    Task theme & Main data or question sources & 1 step & 2--3 & $\geq4$ & Total \\
    \midrule
    Three-dimensional Earth structure & LITHO1.0, SEMUCB-WM1, USGS earthquake catalog & 85 & 15 & 0 & 100 \\
    Geoscientific big-data analysis & Selected GeoQuestions1089 subset & 40 & 0 & 0 & 40 \\
    Critical minerals & Habitable-Earth questions; MRDS and cross-theme climate data & 44 & 14 & 0 & 58 \\
    Climate change & Habitable-Earth questions; paleoclimate, GISTEMP, SSP, and cross-theme mineral data & 31 & 16 & 1 & 48 \\
    Energy resources & MRDS, GEM, OPSD, World Bank, and population data & 100 & 0 & 0 & 100 \\
    Global ocean carbon cycle & SOCAT, GLODAP, SNAPO-CO$_2$ & 100 & 0 & 0 & 100 \\
    Deep-time biosphere & BioDeepTime & 12 & 7 & 5 & 24 \\
    Carbon capture and storage & Synthetic SCCS-OFPOSDU CCS--MRV data & 12 & 5 & 7 & 24 \\
    Remote sensing & Selected Earth-Bench subset & 0 & 7 & 42 & 49 \\
    Natural-disaster event statistics & EM-DAT International Disaster Database & 0 & 55 & 5 & 60 \\
    \midrule
    \textbf{Total} & & \textbf{424} & \textbf{119} & \textbf{60} & \textbf{603} \\
    \bottomrule
  \end{tabularx}
\end{table*}


\subsection{Composition and Evaluation}
\label{sec:metrics}
The benchmark contains 603 evaluation records: 424 have one reference call, 119 have two or three, and 60 have at least four. Its 179 multistep records constitute 29.68\% of the benchmark. Nine cross-theme questions occur in both critical minerals and climate change, so 603 does not denote 603 distinct questions.

A step denotes one reference tool invocation, rather than one scientific operation or a direct measure of scientific difficulty. Energy-resource and ocean-carbon task wrappers, for example, can perform several operations internally. Chain-length strata therefore describe orchestration demands under the current interfaces. A reference chain is one valid implementation; a different valid order is not automatically incorrect. Domain and chain-length groups describe task backgrounds and interface-level workflow composition, respectively.

Final Answer Score (FAS) lies in $[0,1]$ and retains partial credit. Scoring normalizes multiple-choice answers, matches dates and identifiers, applies numerical tolerances to scalars, scores matches among multiple values or text items, and compares structured fields. Empty answers receive zero. For $N$ records,
\begin{equation}
\overline{FAS} = \frac{1}{N}\sum_{i = 1}^{N}{FAS}_{i}
\label{eq:fas}
\end{equation}
The main experiment uses $N=603$; group scores average their respective records with failures retained. This task-level micro-average differs from domain-balanced macro-averages and full-credit rates. FAS assesses agreement with reference answers. Process analysis additionally requires actual parameters, outputs, and traces to assess operation selection, bindings, and artifact consumption.
